% !TeX root = ../main.tex
\section{Evaluation}
\label{sec:evaluation}

In this section, we evaluate our determinization approach in two complementary ways. We seek to addresses the following questions:
\begin{itemize}
    \item \textbf{RQ1: Expressiveness.} What classes of probabilistic programs can our determinization transformation handle? Can determinization enable exact expected value analysis for programs that prior tools can only bound, time out on, or do not support?
    \item \textbf{RQ2: Sampling.} Does determinization improve sampling-based inference without changing the inference algorithm?
\end{itemize}

All experiments have
been run on a (TODO). Runtime
measurements are averaged over 10 executions.

\subsection{Expressiveness}

First, we study \emph{expressiveness}: whether eliminating $\E$-mode randomness can transform probabilistic programs that are difficult to analyze directly into programs amenable to exact expected value analysis. In particular, we are interested in cases where continuous or otherwise infinite state behavior is removed, yielding a finite state probabilistic model that can be analyzed exactly by Storm -- a probabilistic model checker (CITE). 

Our benchmark suite is drawn primarily from prior work on expected value and moment analysis
(CITE ALL: chakarov 2013/14, kura, moosbrugger, beutner). We modify two existing benchmarks, \texttt{asym-random-walk} and \texttt{add-uniform-counter-large}, and introduce three new programs, \texttt{service-policy}, \texttt{uniform-nested}, and \texttt{workload}, designed to exercise features not well represented in existing suites, such as higher order functions and lists. Otherwise, we preserve the original programs and only translate them into our source language. We focus on programs with continuous sampling inside loops or recursion, where the source program may have an uncountable state space but determinization can expose a finite discrete structure.

We distinguish \emph{full} determinization, which eliminates all sampling, from \emph{partial} determinization, which eliminates at least one but not all samples. Partial determinization may nevertheless yield a finite state model when the remaining randomness is finite valued. 

\begin{table}[H]
    \centering
    \small
    \renewcommand{\arraystretch}{1.15}
    \begin{tabular*}{\linewidth}{@{}l@{\extracolsep{\fill}}c c c c r@{}}
        \toprule
        Benchmark & Det. & Expectation &
        \shortstack{Sampling eliminated\\(source $\rightarrow$ det.)} &
        \shortstack{\# states in\\Markov chain} &
        \shortstack{Runtime (s)\\model / Storm} \\
        \midrule
        \texttt{aggregate-rv}     & $\bullet\!\bullet$ & $\mathbb{E}[x]=1.0$                     & $0\!\times\!\mathbb{R} \rightarrow 0$          & 90    & $0.067 / 0.164$ \\
        \texttt{asym-random-walk} & $\bullet\!\circ$   & $\mathbb{E}[t]=0.4$                     & $1\!\times\!\mathbb{R}^{3} \rightarrow 1$      & 200   & $0.185 / 0.169$ \\
        \texttt{coin\_flip\_unif} & $\bullet\!\circ$   & $\mathbb{E}[\mathit{ret}]=1.0$          & $1\!\times\!\mathbb{R} \rightarrow 1$                                 & 72    & $0.068 / 0.161$ \\
        \texttt{dreckon}          & $\bullet\!\circ$   & $\mathbb{E}[x]=0.0$                     & $1\!\times\!\mathbb{R}^{14} \rightarrow 1$     & 43219 & $175.775 / 2.226$ \\
        \texttt{invpend}          & $\bullet\!\bullet$ & $\mathbb{E}[\mathit{pAng}]=2.4$         & $0\!\times\!\mathbb{R}^{6} \rightarrow 0$      & 1822  & $7.061 / 0.245$ \\
        \texttt{pack}             & $\bullet\!\circ$   & $\mathbb{E}[\mathit{totalWeight}]=11.8$ & $1\!\times\!\mathbb{R}^{3} \rightarrow 1$      & 15192 & $37.207 / 1.033$ \\
        \texttt{retransmit}       & $\bullet\!\circ$   & $\mathbb{E}[\mathit{ret}]=0.2$          & $4\!\times\!\mathbb{R}^{2} \rightarrow 4$      & 182   & $0.164 / 0.183$ \\
        \texttt{service\_policy}  & $\bullet\!\circ$   & $\mathbb{E}[\mathit{ret}]=13.2$         & $4\!\times\!\mathbb{R}^{4} \rightarrow 4$      & 9491  & $27.198 / 0.576$ \\
        \texttt{swap}             & $\bullet\!\bullet$ & $\mathbb{E}[x_N]=5.5$                   & $2\!\times\!\mathbb{R}^{2} \rightarrow 0$      & 105   & $0.107 / 0.168$ \\
        \texttt{uniform-counter}  & $\bullet\!\circ$   & $\mathbb{E}[x]=2.0$                     & $1\!\times\!\mathbb{R}^{2} \rightarrow 1$      & 339   & $0.675 / 0.346$ \\
        \texttt{uniform\_nested}  & $\bullet\!\bullet$ & $\mathbb{E}[d]=0.1$                     & $0\!\times\!\mathbb{R}^{4} \rightarrow 0$      & 31    & $0.040 / 0.166$ \\
        \texttt{workload}         & $\bullet\!\circ$   & $\mathbb{E}[\mathit{ret}]=10.0$         & $1\!\times\!\mathbb{R}^{2} \rightarrow 1$      & 27836 & $27.781 / 1.401$ \\
        \bottomrule
    \end{tabular*}
    \caption{Evaluation of determinization on 12 benchmarks. In the determinization
    column, $\bullet\!\bullet$ denotes full determinization and
    $\bullet\!\circ$ denotes partial determinization. In the state space structure column represented as $|D|\times\mathbb{R}^{k}$, $|D|$ is the number of static discrete sampling sites and $k$ is the number of static continuous sampling sites. Runtime excludes both the initial build and Lean certificate checking.}
    \label{tab:finite-after-determinization}
\end{table}

\paragraph{Results} TODO.

\subsection{Sampling}
We next study the utility of determinization as a preprocessing transformation for \emph{sampling-based inference}: whether applying determinization reduces estimator variance. For each benchmark, we compare sampling from the original program $P$ against sampling from $\det(P)$ using the \emph{same} Monte Carlo sampling method. Because determinization analytically eliminates $\E$-mode randomness, the transformed program produces lower variance estimators without requiring any modification to the underlying inference procedure. We report the empirical estimator variance and variance reduction factor on all programs.

We introduce three new benchmarks, \texttt{hmm}, \texttt{load}, and \texttt{branching}, and use modified versions of two existing programs, \texttt{neals-funnel} and \texttt{growing-walk}. The suite includes both full and partial determinization and covers reductions from continuous state to finite, countably infinite, and lower-dimensional continuous state. 

\begin{table}[H]
    \centering
    \small
    \renewcommand{\arraystretch}{1.15}
    \begin{tabular*}{\linewidth}{@{}l@{\extracolsep{\fill}}c c c c@{}}
        \toprule
        Benchmark &
        \shortstack{State-space structure\\(source $\rightarrow$ det.)} &
        $\widehat{\operatorname{Var}}(P)$ &
        $\widehat{\operatorname{Var}}(\det P)$ & VRF \\
        \midrule
        \texttt{2d-walk}             & $1\!\times\!\mathbb{R} \rightarrow 1$                       & 167.16 & 125.37 & $1.33\times$ \\
        \texttt{bimodal}             & $1\!\times\!\mathbb{R}^{2} \rightarrow 1$                   & 4.53   & 0.50   & $9.02\times$ \\
        \texttt{branching}           & $1\!\times\!\mathbb{R} \rightarrow 1$                       & 1.66   & 0.39   & $4.25\times$ \\
        \texttt{continuous-sampling} & $1\!\times\!\mathbb{R}^{3} \rightarrow 1\!\times\!\mathbb{R}$ & 38.01 & 0.00 & $\infty$ \\
        \texttt{ckd-epi}             & $0\!\times\!\mathbb{R}^{8} \rightarrow 0\!\times\!\mathbb{R}^{6}$ & 0.00 & 0.00 & $1.00\times$ \\
        \texttt{figure1}             & $3\!\times\!\mathbb{R}^{3} \rightarrow 3\!\times\!\mathbb{R}$ & 0.21 & 0.21 & $1.00\times$ \\
        \texttt{growing-walk}        & $1\!\times\!\mathbb{R} \rightarrow 1$                       & 6.12   & 5.62   & $1.09\times$ \\
        \texttt{hmm}                 & $2\!\times\!\mathbb{R} \rightarrow 2$                       & 12.30  & 0.11   & $109.73\times$ \\
        \texttt{load}                & $0\!\times\!\mathbb{R}^{5} \rightarrow 0\!\times\!\mathbb{R}$ & 58.35 & 1.44 & $40.58\times$ \\
        \texttt{neals-funnel}        & $0\!\times\!\mathbb{R}^{11} \rightarrow 0\!\times\!\mathbb{R}$ & 1.46 & 0.00 & $\infty$ \\
        \texttt{random-walk-1d}      & $1\!\times\!\mathbb{R} \rightarrow 1$                       & 1.73   & 0.90   & $1.93\times$ \\
        \texttt{vasicek}             & $0\!\times\!\mathbb{R} \rightarrow 0$                       & 0.85   & 0.00   & $\infty$ \\
        \bottomrule
    \end{tabular*}
    \caption{Variance reduction on the sampling benchmark suite using 100,000
    attempted runs per program. Empirical variances are computed among returned
    values. State-space structures use the same static sampling-site notation as
    Table~\ref{tab:finite-after-determinization}. The variance-reduction factor
    is $\mathrm{VRF}=\widehat{\operatorname{Var}}(P) /
    \widehat{\operatorname{Var}}(\det P)$.}
    \label{tab:infinite-after-determinization}
\end{table}

\paragraph{Results} TODO.
